\paragraph{Phát hiện}

\def\TP{\relax\ifmmode\mathrm{TP}\else TP\fi}
\def\FP{\relax\ifmmode\mathrm{FP}\else FP\fi}
\def\FN{\relax\ifmmode\mathrm{FN}\else FN\fi}

Để đánh giá chất lượng của mô hình phát hiện đối tượng, các độ đo được sử
dụng bao gồm Precision, Recall, F1-score, mAP@50 và mAP@50--95. Các độ
đo này phản ánh khả năng phát hiện, phân loại và định vị đối tượng của
mô hình trên tập dữ liệu đánh giá.

Precision đánh giá tỷ lệ các đối tượng được mô hình dự đoán là đúng trên
tổng số các đối tượng được dự đoán. Chỉ số này phản ánh mức độ tin cậy
của các dự đoán do mô hình tạo ra và được tính như sau:
\begin{equation}
\mathrm{Precision}
= \frac{\TP}{\TP + \FP},
\end{equation}
trong đó, \TP\ (True Positives) là số lượng đối tượng được phát hiện
đúng và \FP\ (False Positives) là số lượng đối tượng được dự đoán nhưng
không tồn tại trong thực tế.

Recall đánh giá khả năng phát hiện đầy đủ các đối tượng có trong ảnh.
Chỉ số này được tính bằng tỷ lệ số đối tượng được phát hiện đúng trên
tổng số đối tượng thực tế:
\begin{equation}
\mathrm{Recall}
= \frac{\TP}{\TP + \FN},
\end{equation}
trong đó, \FN\ (False Negatives) là số lượng đối tượng thực tế nhưng bị
mô hình bỏ sót. Recall cao cho thấy mô hình có khả năng phát hiện được
phần lớn các đối tượng trong dữ liệu.

F1-score là trung bình điều hòa của Precision và Recall, được sử dụng để
đánh giá sự cân bằng giữa độ chính xác và khả năng phát hiện của mô
hình:
\begin{equation}
\mathrm{F1}
= \frac{2 \times \mathrm{Precision} \times \mathrm{Recall}}
{\mathrm{Precision} + \mathrm{Recall}}.
\end{equation}
Giá trị F1-score càng cao cho thấy mô hình đạt được sự cân bằng tốt giữa
Precision và Recall.

Trong bài toán phát hiện đối tượng, đường cong Precision--Recall (PR)
được xây dựng bằng cách thay đổi ngưỡng confidence của mô hình. Từ
đường cong này, diện tích dưới đường cong được sử dụng để tính Average
Precision (AP):
\begin{equation}
  \mathrm{AP} = \int_0^1 P(R)\, dR,
\end{equation}
trong đó $P(R)$ là giá trị Precision tương ứng với Recall $R$. AP phản
ánh đồng thời khả năng phát hiện và phân loại đối tượng của mô hình trên
một lớp cụ thể.

Đối với bài toán nhiều lớp, chỉ số Mean Average Precision (mAP) được
tính bằng giá trị trung bình của AP trên tất cả các lớp:
\begin{equation}
  \mathrm{mAP} = \frac{1}{N}\sum_{i=1}^{N} \mathrm{AP}_i,
\end{equation}
trong đó $N$ là số lớp đối tượng và $AP_i$ là Average Precision của lớp
thứ $i$.

Chỉ số mAP@50 được tính tại ngưỡng IoU bằng 0.5, nghĩa là một đối tượng
được xem là phát hiện đúng nếu độ chồng lấp giữa hộp dự đoán và hộp thực
tế đạt ít nhất 50%. Đây là độ đo được sử dụng phổ biến để đánh giá khả
năng phát hiện tổng quát của mô hình.

Để đánh giá chặt chẽ hơn chất lượng định vị, chỉ số mAP@50--95 được sử
dụng. Độ đo này được tính bằng giá trị trung bình của AP tại các ngưỡng
IoU từ 0.5 đến 0.95 với bước nhảy 0.05:

\begin{equation}
\mathrm{mAP}@50\text{--}95
= \frac{1}{10}
\sum_{\mathrm{IoU}\in{0.50,0.55,\ldots,0.95}}
\mathrm{AP}_{\mathrm{IoU}}.
\end{equation}

Do được tính trên nhiều ngưỡng IoU khác nhau, mAP@50--95 phản ánh toàn
diện hơn khả năng phát hiện và định vị đối tượng của mô hình.
